\documentclass[10pt,journal,compsoc]{IEEEtran}

% Basic packages
\usepackage[utf8]{inputenc}
\usepackage[T1]{fontenc}
\usepackage{amsmath,amssymb,amsfonts}
\usepackage{algorithm}
\usepackage{algpseudocode}
\usepackage{graphicx}
\usepackage{xcolor}
\usepackage{booktabs}
\usepackage{multirow}
\usepackage{subcaption}
\usepackage{hyperref}
\usepackage{cleveref}
\usepackage{balance}

% Theorem environments
\usepackage{amsthm}
\newtheorem{theorem}{Theorem}
\newtheorem{lemma}{Lemma}
\newtheorem{definition}{Definition}
\newtheorem{property}{Property}

% Custom commands
\newcommand{\systemname}{AetherGuard}
\newcommand{\agentset}{\mathcal{A}}
\newcommand{\threatset}{\mathcal{T}}
\newcommand{\actionset}{\mathcal{X}}
\newcommand{\decision}{\mathcal{D}}

% Color definitions
\definecolor{darkblue}{RGB}{0,51,102}
\definecolor{darkgreen}{RGB}{0,102,51}
\definecolor{darkred}{RGB}{153,0,0}

\hypersetup{
    colorlinks=true,
    linkcolor=darkblue,
    citecolor=darkblue,
    urlcolor=darkblue
}

\begin{document}

% Title
\title{\systemname: A Multi-Agent AI Architecture for Distributed Server Security}

% Authors (placeholder)
\author{
    \IEEEauthorblockN{Anonymous Authors}
    \IEEEauthorblockA{
        Institution Name\\
        Email: anonymous@institution.edu
    }
}

% Make title
\maketitle

% Abstract
\begin{abstract}
% To be written by writer-abstract-intro agent
% 150-250 words covering:
% - Problem statement
% - Core innovation (7-layer AI Agent architecture)
% - Key technical contributions
% - Experimental results highlights
% - Practical significance
\textbf{TODO:} Abstract to be written by writer-abstract-intro agent.
\end{abstract}

% Keywords
\begin{IEEEkeywords}
Distributed Security, Multi-Agent Systems, Artificial Intelligence, Intrusion Detection, Real-time Response, Server Protection
\end{IEEEkeywords}

% Section 1: Introduction
\section{Introduction}
\label{sec:introduction}
% To be written by writer-abstract-intro agent
\textbf{TODO:} Introduction to be written by writer-abstract-intro agent.

\subsection{Research Motivation}
\label{subsec:motivation}

\subsection{Research Challenges}
\label{subsec:challenges}

\subsection{Contributions}
\label{subsec:contributions}
The main contributions of this paper are:
\begin{enumerate}
    \item \textbf{TODO:} Contribution 1 to be written by innovation-analyst agent.
    \item \textbf{TODO:} Contribution 2 to be written by innovation-analyst agent.
    \item \textbf{TODO:} Contribution 3 to be written by innovation-analyst agent.
    \item \textbf{TODO:} Contribution 4 to be written by innovation-analyst agent.
    \item \textbf{TODO:} Contribution 5 to be written by innovation-analyst agent.
\end{enumerate}

\subsection{Paper Organization}
\label{subsec:organization}

% Section 2: Background and Related Work
\section{Background and Related Work}
\label{sec:related-work}
% To be written by writer-background-related agent
\textbf{TODO:} Background and Related Work to be written by writer-background-related agent.

\subsection{Distributed Intrusion Detection Systems}
\label{subsec:distributed-ids}

\subsection{Agent-based Security Architectures}
\label{subsec:agent-security}

\subsection{AI-driven Security Systems}
\label{subsec:ai-security}

\subsection{Multi-Agent Collaboration Mechanisms}
\label{subsec:multi-agent}

\subsection{Limitations of Existing Approaches}
\label{subsec:limitations}

% Section 3: System Architecture
\section{System Architecture}
\label{sec:architecture}
% To be written by writer-system-design agent
\textbf{TODO:} System Architecture to be written by writer-system-design agent.

\subsection{System Overview}
\label{subsec:overview}

\subsection{Control Plane Design}
\label{subsec:control-plane}

\subsection{Distributed Agent Design}
\label{subsec:agent-design}

\subsection{Communication Architecture}
\label{subsec:communication}

% Section 4: Multi-Agent AI Architecture (Core Innovation)
\section{Multi-Agent AI Architecture}
\label{sec:ai-architecture}
% To be written by writer-system-design agent
\textbf{TODO:} Multi-Agent AI Architecture to be written by writer-system-design agent.

\subsection{Seven-Layer Architecture}
\label{subsec:seven-layer}

\subsubsection{Layer 1: Collaboration Layer}
\label{subsubsec:layer1}

\subsubsection{Layer 2: Role Playing Layer}
\label{subsubsec:layer2}

\subsubsection{Layer 3: Scene Traversal Layer}
\label{subsubsec:layer3}

\subsubsection{Layer 4: Capability Layer}
\label{subsubsec:layer4}

\subsubsection{Layer 5: Debate Layer}
\label{subsubsec:layer5}

\subsubsection{Layer 6: Heterogeneous Model Layer}
\label{subsubsec:layer6}

\subsubsection{Layer 7: Synthesis Layer}
\label{subsubsec:layer7}

\subsection{Agent Roles}
\label{subsec:agent-roles}

\subsection{Debate and Arbitration Mechanism}
\label{subsec:debate}

\subsection{Self-Evolution Capability}
\label{subsec:self-evolution}

% Section 5: Algorithm Design
\section{Algorithm Design}
\label{sec:algorithm}
% To be written by algorithm-designer and mathematical-formalist agents
\textbf{TODO:} Algorithm Design to be written by algorithm-designer and mathematical-formalist agents.

\subsection{Formal Definitions}
\label{subsec:definitions}

\begin{definition}[Multi-Agent System]
\label{def:mas}
A multi-agent system $\mathcal{M}$ is defined as a tuple:
\begin{equation}
\mathcal{M} = (\agentset, \mathcal{S}, \mathcal{O}, \mathcal{P}, \mathcal{R})
\end{equation}
where $\agentset$ is the set of agents, $\mathcal{S}$ is the state space, $\mathcal{O}$ is the observation space, $\mathcal{P}$ is the transition probability, and $\mathcal{R}$ is the reward function.
\end{definition}

\subsection{Multi-Agent Decision Algorithm}
\label{subsec:decision-algorithm}

\begin{algorithm}
\caption{Multi-Agent Decision Algorithm}
\label{alg:decision}
\begin{algorithmic}[1]
\Require Threat observation $o \in \mathcal{O}$, Agent set $\agentset$
\Ensure Decision $\decision$
\State \textbf{TODO:} Algorithm to be written by algorithm-designer agent.
\end{algorithmic}
\end{algorithm}

\subsection{Threat Detection Algorithm}
\label{subsec:detection-algorithm}

\subsection{Command Signing and Verification}
\label{subsec:signing}

\subsection{Complexity Analysis}
\label{subsec:complexity}

% Section 6: Security Analysis
\section{Security Analysis}
\label{sec:security}
% To be written by security-analyst agent
\textbf{TODO:} Security Analysis to be written by security-analyst agent.

\subsection{Threat Model}
\label{subsec:threat-model}

\subsection{Attacker Model}
\label{subsec:attacker-model}

\subsection{Security Properties}
\label{subsec:security-properties}

\begin{property}[Authentication]
\label{prop:authentication}
\textbf{TODO:} To be written by security-analyst agent.
\end{property}

\begin{property}[Integrity]
\label{prop:integrity}
\textbf{TODO:} To be written by security-analyst agent.
\end{property}

\subsection{Formal Security Proofs}
\label{subsec:proofs}

\begin{theorem}[Security Guarantee]
\label{thm:security}
\textbf{TODO:} To be written by security-analyst agent.
\end{theorem}

% Section 7: Implementation
\section{Implementation}
\label{sec:implementation}
% To be written by writer-implementation agent
\textbf{TODO:} Implementation to be written by writer-implementation agent.

\subsection{Technology Stack}
\label{subsec:tech-stack}

\subsection{Key Implementation Details}
\label{subsec:implementation-details}

\subsection{Deployment Architecture}
\label{subsec:deployment}

% Section 8: Evaluation
\section{Evaluation}
\label{sec:evaluation}
% To be written by writer-evaluation agent
\textbf{TODO:} Evaluation to be written by writer-evaluation agent.

\subsection{Experimental Setup}
\label{subsec:setup}

\subsection{Performance Metrics}
\label{subsec:metrics}

\subsection{System Performance}
\label{subsec:system-performance}

\subsubsection{Latency Analysis}
\label{subsubsec:latency}

\subsubsection{Throughput Analysis}
\label{subsubsec:throughput}

\subsubsection{Scalability Analysis}
\label{subsubsec:scalability}

\subsection{AI Agent Performance}
\label{subsec:ai-performance}

\subsection{Comparison with Existing Systems}
\label{subsec:comparison}

\begin{table}[t]
\caption{Comparison with Existing Systems}
\label{tab:comparison}
\centering
\begin{tabular}{lcccc}
\toprule
\textbf{System} & \textbf{Latency} & \textbf{Accuracy} & \textbf{Scalability} & \textbf{AI-driven} \\
\midrule
\textbf{TODO:} & - & - & - & - \\
\systemname & - & - & - & \checkmark \\
\bottomrule
\end{tabular}
\end{table}

\subsection{Ablation Study}
\label{subsec:ablation}

\subsection{Real-world Deployment Results}
\label{subsec:real-world}

% Section 9: Discussion
\section{Discussion}
\label{sec:discussion}

\subsection{Limitations}
\label{subsec:limitations}

\subsection{Future Work}
\label{subsec:future}

% Section 10: Conclusion
\section{Conclusion}
\label{sec:conclusion}
% To be written by writer-conclusion-refs agent
\textbf{TODO:} Conclusion to be written by writer-conclusion-refs agent.

% References
\bibliographystyle{IEEEtran}
% \bibliography{references}  % Uncomment when references.bib is ready

% Placeholder references
\begin{thebibliography}{10}
\bibitem{example1}
\textbf{TODO:} References to be written by writer-conclusion-refs agent.
\end{thebibliography}

% Appendix
\appendix

\section{Additional Mathematical Proofs}
\label{app:proofs}
\textbf{TODO:} To be written by mathematical-formalist agent.

\section{Supplementary Experimental Results}
\label{app:experiments}
\textbf{TODO:} To be written by writer-evaluation agent.

\section{Implementation Details}
\label{app:implementation}
\textbf{TODO:} To be written by writer-implementation agent.

\section{Dataset Description}
\label{app:dataset}
\textbf{TODO:} To be written by writer-evaluation agent.

\end{document}
